\documentclass[12pt]{article}

\input{../../_assets/preambulo.tex}


\begin{document}

    % 1. Foto de fondo
    % 2. Título
    % 3. Encabezado Izquierdo
    % 4. Color de fondo
    % 5. Coord x del titulo
    % 6. Coord y del titulo
    % 7. Fecha

    
    \input{../../_assets/portada}
    \portadaExamen{ffccA4.jpg}{Ecuaciones\\Diferenciales II\\Examen XXIII}{Ecuaciones Diferenciales II. Examen XXIII}{MidnightBlue}{-8}{28}{2026}{}

    \begin{description}
        \item[Asignatura] Ecuaciones Diferenciales II.
        \item[Curso Académico] 2023/2024.
        \item[Grado] Doble Grado en Ingeniería Informática y Matemáticas.
        \item[Grupo] Único.
        \item[Profesor] Rafael Ortega Ríos.
        \item[Descripción] Convocatoria Extraordinaria.
        \item[Fecha] 9 de Julio de 2024.
        % \item[Duración] ---.
    
    \end{description}
    \newpage


    % ------------------------------------
    
    \begin{ejercicio}
        Decide en cada caso si la sucesión de funciones $\{f_n\}_{n \geq 1}$, $f_n : \left[0, 1\right] \to \bb{R}$, es equicontinua.
        \begin{enumerate}[label=\alph*)]
            \item $f_n(t) = t^n$
            \item $f_n(t) = \sqrt{t + \frac{1}{n}}$
        \end{enumerate}
    \end{ejercicio}

    \begin{ejercicio}
        Encuentra todas las funciones continuas $x, y : \bb{R} \to \bb{R}$ que cumplen
        \[
            x(t) = 3 + \int_{0}^{t} y(s)ds, \quad y(t) = -4 \int_{0}^{t} x(s)ds.
        \]
    \end{ejercicio}

    \begin{ejercicio}
        Se supone que $x(t)$ es la solución maximal del problema
        \[
            \dot{x} = 2x + \frac{\cos t}{x^2}, \quad x(0) = 2.
        \]
        \begin{enumerate}[label=\alph*)]
            \item Demuestra que $\dot{x}(t) > 0$ para cada $t \in \left[0, \omega\right[$.
            \item $\omega = +\infty$.
        \end{enumerate}
    \end{ejercicio}

    \begin{ejercicio}
        Para cada $n \geq 1$ se denota por $x_n(t)$ a la solución maximal del problema
        \[
            \dot{x} = x^5, \quad x(0) = 1 - \frac{1}{n}.
        \]
        Calcula, si existe, $\lim\limits_{n \to \infty} x_n(t)$ para cada $t \in \left[0, 1\right[$. ¿Hay convergencia uniforme en $\left[0, 1\right[$?
    \end{ejercicio}

    \begin{ejercicio}
        Encuentra una función continua $f : \bb{R} \times \bb{R} \to \bb{R}$ de manera que el problema de valores iniciales
        \[
            x' = f(t, x), \quad x(2) = 3
        \]
        tenga más de una solución maximal.
    \end{ejercicio}

    \begin{ejercicio}
        Dado un número $\gamma \in \left]0, \nicefrac{1}{2}\right[$ y una función continua y acotada dada por $b : \bb{R} \to \bb{R}$, se considera la ecuación
        \[
            x(t) = \gamma(x(t - 1) + x(t + 1)) + b(t).
        \]
        Una solución $x : \bb{R} \to \bb{R}$ es una función continua y acotada que cumple esta identidad para cada $t \in \bb{R}$.
        \begin{enumerate}[label=\alph*)]
            \item Demuestra que existe a lo sumo una solución.
            \item Demuestra que existe al menos una solución.
        \end{enumerate}
        \begin{observacion}
            Considerar $x_{n+1}(t) = \gamma(x_n(t - 1) + x_n(t + 1)) + b(t)$.
        \end{observacion}
    \end{ejercicio}

    \newpage
    \setcounter{ejercicio}{0} % Reiniciar contador de ejercicios
    \noindent
    % \textbf{Solución.}

\end{document}